\subsection{Ausblick}
\label{subsec:ausblick}

Kurzfristig sind die roten Tooling-, Desktopprofil- und Windows-CI-Pfade auf
einem festen Commit zu korrigieren. Compose-Start mit Health und Readiness,
Release Validation und Branch Protection sollen die technische Evidenz
ergänzen, ohne Produktreife zu beanspruchen.

Mittelfristig benötigen konkrete Produkte Signierung, Notarisierung und einen
Updater. Auch Credentials, Deployment, Authentifizierung und RBAC bleiben
produktspezifisch. Für die Baseline bieten sich drei Schritte an: gemeinsame
Contracts automatisieren, Capabilities strenger wählen und generierte Projekte
kontrolliert synchronisieren.

Langfristig sollte eine projektübergreifende Vergleichsstudie Onboarding,
ersten Build, Setup-Fehler und Pflegeaufwand messen. Reale Kundenprojekte
könnten die Profil- und Eignungskriterien prüfen.
